\section{Alternatives and Research Position}
Touch and haptic authentication have a substantial history. TouchIn combines
knowledge of a gesture with behavioral and physiological features~\cite{touchin}.
Recent tele-robotic work uses force-feedback dynamics for user
identification~\cite{haptic}. OpenATT is neither a biometric classifier nor a
competitor to their user-identification accuracy; it has no participant dataset.
Its pattern is public and contains only eight bits, so it is not an identity
secret.

WebAuthn separately specifies user-presence and user-verification checks and
relying-party validation~\cite{webauthn}. An arbitrary force threshold or
controller signature does not instantiate that ceremony. A trusted authorization
flow can supply permission; a tactile event can supply a distinct reported
observation. Combining them must preserve their meanings.

RATS separates evidence from appraisal~\cite{rats}. Analog sensor-injection
research such as WALNUT demonstrates that digital processing can begin with
adversarial physical inputs~\cite{walnut}; a signature downstream cannot establish
the integrity of that upstream physical measurement. CaMeL and Fides enforce
control/data-flow policies around agents~\cite{camel,fides}. OpenATT uses a
trusted action gate but does not implement their planners, prove their properties
or reproduce their agent evaluation.

Accordingly, signed controller reports and exact grants are engineering
foundations. The prospective scientific question is whether a common evidence
contract can preserve non-equivalent guarantees across changing perceptual
contexts and acquire adequate support at lower cost. The current work supplies
an inspectable tactile boundary and limitations, not that general result.

